\chapter{Methodology}
\label{ch:methodology}

\section{Introduction}
\label{sec:methodology-intro}

\section{Research Design}
\label{sec:research-design}

\section{Proposed Approach}
\label{sec:approach}

Coroutines' compilation has 3 points of interest:

\begin{enumerate}
    \item Transforming suspend functions into state-machines.
    \item Adding a continuation parameter to suspend functions' signatures.
    \item Providing compiler intrinsics for coroutines.
\end{enumerate}

In order to achieve our goals we need to do two things:

\begin{enumerate}
    \item Disable transformation of suspend functions into state-machines.
    \item Implement compiler intrinsics with the use of Wasm continuations.
\end{enumerate}

\subsection{Main Challenges}
\label{subsec:challenges}

The main challenges are:

\begin{itemize}
    \item Figuring out a compilation scheme with Wasm continuations that doesn't rely on state-machine transformation;
    \item Seamlessly integrating the changes into the Kotlin compiler without breaking the compiler's functionality;
    \item Ensuring that the new implementation for suspend functions and coroutine intrinsics doesn't change their behavior (to the extent that is necessary for achieving the ultimate goal).
\end{itemize}

The last point can be tested with the existing testing suite of the Kotlin compiler.

\section{Algorithms}
\label{sec:algorithms}

\section{Data Collection and Processing}
\label{sec:data}

\section{Experimental Setup}
\label{sec:experimental-setup}

\section{Evaluation Metrics}
\label{sec:metrics}

\section{Summary}
\label{sec:methodology-summary}
